\begin{appendix}

\chapter{Implementation and Reproducibility}
\label{app:implementation}

All computations are done in Python. The estimators and the inference procedure form a
small package, \texttt{rankci}; Jupyter notebooks and two scripts call it to run the
Monte Carlo study and the applications. Below I map the method onto the code
(\S\ref{app:impl-code}), record the implementation choices the main text leaves open
(\S\ref{app:impl-details}), list the data (\S\ref{app:impl-data}), the software and the
tests (\S\ref{app:impl-env}), state which file produces each result
(\S\ref{app:impl-repro}), and list the known limitations
(\S\ref{app:impl-limits}).

\section{From the method to the code}
\label{app:impl-code}

The project directory has six parts. All of them except \texttt{data/} and
\texttt{thesis/} are under version control.
\begin{itemize}[nosep,leftmargin=*]
  \item \texttt{src/rankci/} is the package. Its subpackage \texttt{core/} holds the
    inference engines, which take a plain $(n\times p)$ array and know nothing about
    where the data came from. \texttt{data/} holds one loader per data source, and
    \texttt{sim.py} the Monte Carlo design.
  \item \texttt{tests/} holds the unit tests (\S\ref{app:impl-tests}).
  \item \texttt{notebooks/} holds the Monte Carlo study
    (\texttt{01\_toydataset.ipynb}), the dependence analysis of
    Appendix~\ref{app:spf-dependence} (\texttt{02\_primary\_analysis.ipynb}), and one
    folder per application (\texttt{philly/}, \texttt{ecb/}, \texttt{gjp/},
    \texttt{mcomp/}).
  \item \texttt{scripts/} holds \texttt{make\_thesis\_figures.py}, which regenerates
    Figures~\ref{fig:mc-crossover}, \ref{fig:gjp-caterpillar}
    and~\ref{fig:mcomp-caterpillar}, and \texttt{make\_thesis\_tables.py}, which
    regenerates Tables~\ref{tab:mc-coverage}, \ref{tab:mc-crossover},
    \ref{tab:spf-indicators}, \ref{tab:ngdp}
    and~\ref{tab:threeregimes} and every table of Appendix~\ref{app:additional}, and
    prints the euro-area results and the NGDP $\tau$-best sets.
  \item \texttt{thesis/} holds this document and nothing else: its LaTeX sources and
    the figures and tables that the notebooks and scripts write into it.
  \item \texttt{data/} holds the raw data (\S\ref{app:impl-data}).
\end{itemize}
Table~\ref{tab:impl-code} gives the function that implements each step of the
pipeline of \S\ref{sec:pipeline}. A complete run takes a loss panel and returns the
rank confidence sets:
\begin{verbatim}
from rankci import rank_ci_stepwise_simulation_pairwise
out = rank_ci_stepwise_simulation_pairwise(
    X, alpha=0.1, B=20000, seed=42, covariance="omega")
out["rank_ci"]   # (p, 2): lower and upper rank of each column
\end{verbatim}
Here \texttt{X} is the $(n\times p)$ loss panel with \texttt{NaN} for missing entries,
and \texttt{covariance="mds"} selects the MDS reconstruction instead. The returned
dictionary also holds the estimated means, the covariance the draws were taken from,
the confirmed pairs, the number of stepdown rounds, and two projection diagnostics: the
distance $\Frob{A-A^{+}}$ between the raw matrix and its projection, and the smallest
eigenvalue of the raw matrix.

\tabImplCode

\section{Implementation choices}
\label{app:impl-details}

\paragraph{Input and conventions.} A loss panel is an $(n\times p)$ floating-point
array whose rows are ordered in time and whose missing entries are \texttt{NaN}. Rank~1
is the smallest loss. The point estimate $\thetahat_j$ is the mean of column $j$ over
its own observed rows, whereas the contrast $\Deltahat_{j,k}$ uses only the overlap
$\Tjk$ (\S\ref{sec:setup}). On an unbalanced panel the order of the $\thetahat_j$ can
therefore disagree with the signs of the $\Deltahat_{j,k}$; forecaster~65 in
\S\ref{sec:apps-spf} is an example.

\paragraph{Untestable pairs.} A pair with fewer than two common observations (the
default \texttt{min\_overlap=2}), or with a zero standard error, is never tested. It
can never be confirmed, so it can only widen the rank sets of its two units.

\paragraph{Lags on an unbalanced panel.} \texttt{compute\_pairwise} extracts the rows
on which both $j$ and $k$ are observed and concatenates them before evaluating
\eqref{eq:nw}. Lag $\tau$ therefore means $\tau$ joint observations apart, not $\tau$
rows apart. The Andrews bandwidth is computed from the same concatenated series:
$\hat\rho$ is the least-squares slope of $d_t$ on $d_{t-1}$ with an intercept, clamped
to $\pm0.97$, and $L=0$ for fewer than three observations. An explicit bandwidth
overrides the plug-in. Winsorization of the difference series is implemented
(\texttt{winsor\_pct}) but not used for any reported result.

The direct estimator keeps each difference series on the full row grid, with
\texttt{NaN} off its support, and sums the cross-products in \eqref{eq:omega-entry} over
the paired index set $T_{ab}(h)$, with lags counted in rows. Its bandwidths $L_a$ are
the plug-in values of the concatenated series, and
$L_{ab}=\min\{L_a,L_b,n_{ab}-1\}$. The two lag conventions agree when every pair's
overlap is a contiguous block of rows. That holds on a balanced panel and under the
staggered entry and exit of the Monte Carlo design. When an overlap has gaps, they
differ by the autocovariance terms that straddle a gap, so the
diagonal-consistency identity $\Omegahat_{aa}=(\seNW{a})^2$ of
\S\ref{sec:omega-nesting} is exact for contiguous overlaps only. On the NGDP panel of
Table~\ref{tab:ngdp}, whose overlaps have gaps, $\sqrt{\Omegahat_{aa}}$ lies between
$0.95$ and $1.12$ times $\seNW{a}$ before projection.

\paragraph{Studentization.} As described in \S\ref{sec:seam}, the MDS route
studentizes both the observed and the simulated contrasts by $\seNW{a}$ from
\texttt{compute\_pairwise}, and the direct route studentizes both by
$\sqrt{\Omegahat^{+}_{aa}}$, the diagonal of the projected matrix. The direct route's
draws are therefore exactly unit-variance even when the projection is active, but its
observed statistics are then divided by a standard error that differs from $\seNW{a}$.
On the NGDP panel the projection is nearly inert (its Frobenius distance is below
$0.1\%$ of $\Frob{\Omegahat}$), so $\sqrt{\Omegahat^{+}_{aa}}$ stays between $0.95$ and
$1.13$ times $\seNW{a}$, as before projection.

\paragraph{Entries of $\Omegahat$.} Each difference series is demeaned on its own
support. Each Bartlett-weighted cross sum is divided by $n_a n_b$, the product of the
two support sizes, as in \eqref{eq:omega-entry}; on the diagonal this is $n_a^2$.
Off-diagonal entries whose pair-of-pairs overlap $n_{ab}$ is below
\texttt{min\_overlap} are set to zero, while diagonal entries always use the full
support $T_a$. The assembled matrix is symmetrized before the projection.

\paragraph{Nearest-PSD projection.} Both routes use the same function. It
symmetrizes the matrix, computes its eigendecomposition with
\texttt{numpy.linalg.eigh}, sets the negative eigenvalues to zero, reassembles the
matrix and symmetrizes it again. The result is the Frobenius-nearest PSD matrix
\citep{higham1988}.

\paragraph{Simulation and stepdown.} The Gaussian draws come from
\texttt{numpy.random.\allowbreak Generator.\allowbreak multivariate\_normal}. It factors the covariance by a
singular value decomposition and therefore accepts singular matrices such as $\Sh$,
for which $\Sh\one=0$. The $B$ draws are taken once and reused in every round. Each
unordered pair $a=(j,k)$ gives two one-sided hypotheses, $(j,k)$ and $(k,j)$; the
second uses the observed statistic and the draws with reversed sign. In each round the
critical value is the empirical $(1-\alpha)$ quantile of the row-wise maximum over the
active hypotheses (\texttt{numpy.quantile} with linear interpolation), and a
hypothesis is rejected when its statistic strictly exceeds the critical value.

\paragraph{Block bootstrap benchmark.} \texttt{rank\_ci\_stepwise\_pairwise}
implements the circular block bootstrap of \S\ref{par:mbb}
(\texttt{core/block\_bootstrap.py}). Block starts are uniform on $\{1,\dots,n\}$, each
block covers $\ell$ consecutive rows and continues from row $n$ to row $1$, and
$\lceil n/\ell\rceil$ blocks are concatenated and truncated to $n$ rows. Every row is
therefore drawn equally often in expectation; without the wrap-around the first and
last $\ell-1$ rows would be drawn less often, which shifts the bootstrap contrasts away
from the observed ones. The block length defaults to $\ell=L+1$ with $L$ the median
Andrews bandwidth over the pairwise difference series (\texttt{default\_block\_length});
a user-supplied $L$ gives $\ell=L+1$ directly. The $B$ resamples are drawn once
(\texttt{block\_bootstrap\_draws}) and passed to the same stepdown function as the
Gaussian draws. A pair without overlap in a resample does not enter that replication's
maximum. The marginal and $\tau$-best bootstrap variants draw their resamples in the
same way.

\paragraph{Marginal and $\tau$-best sets.} Marginal rank sets are computed in a single
step, with one critical value per unit from the $2(p-1)$ hypotheses that involve it.
The direct $\tau$-best procedure enumerates the exemption sets $K$ with
\texttt{itertools.combinations}, so its cost grows with $\binom{p}{\tau-1}$. Like the
stepdown, it takes the critical value of every round from one set of draws, Gaussian or
bootstrap.

\paragraph{Random numbers.} Every stochastic step draws from a
\texttt{numpy.random.Generator} created from an explicit integer seed. The
applications use seed $42$. In the Monte Carlo study, replication $r$ of design $i$,
with $i$ counted within each study (the eleven-design grid and the six-design crossover
sweep), simulates its panel with seed $10^6 i+r$ and its Gaussian draws, as well as its block
bootstrap resamples, with that seed plus $5\times10^8$. The MDS and direct routes and
the block bootstrap therefore see identical panels, so the comparison of
\S\ref{sec:mc-design} is paired, and the draws are independent of the data.

Results are reproducible on a given machine and library version. On another platform
the linear-algebra library can return a different but equally valid factorization, so
the same seed can give different draws. Critical values then agree up to Monte Carlo
error, and a contrast within that error of the critical value can change its decision.

\section{Data}
\label{app:impl-data}

The data are not part of the code. All are public and available from the sources
below. The loaders expect them under \texttt{data/} with the file names given here and
read them unchanged.
\begin{itemize}[leftmargin=*]
  \item \emph{U.S.\ SPF} (\S\ref{sec:apps-spf}, Appendix~\ref{app:spf-dependence}).
    The individual forecasts of the Survey of Professional Forecasters of the Federal
    Reserve Bank of Philadelphia \citep{croushore2019spf} are in
    \texttt{philly/SPFmicrodata.xlsx}, sheets \texttt{NGDP}, \texttt{RGDP} and
    \texttt{UNEMP}. The realizations come from the Bank's Real-Time Data Set for
    Macroeconomists: \texttt{philly/NOUTPUTQvQd.xlsx} and
    \texttt{philly/ROUTPUTQvQd.xlsx} (quarterly vintages of nominal and real output)
    and \texttt{philly/rucQvMd.xlsx} (monthly vintages of the unemployment rate,
    averaged to quarters). The advance estimate of a quarter is read from the vintage
    dated one quarter later. \texttt{data.philly.compute\_error\_panel} builds the
    loss panel at SPF horizon index $3$, and \texttt{select\_top\_forecasters} keeps
    the forecasters with the most observations among those with at least $20$.
  \item \emph{Euro-area SPF} (\S\ref{sec:apps-ecb}). The individual replies to the
    ECB Survey of Professional Forecasters are one CSV file per survey round,
    \texttt{ecb/individual\_forecasts/YYYYQq.csv} (110 rounds, 1999Q1--2026Q2), read
    by \texttt{data.ecb.load\_ecb\_spf}. The realizations are Eurostat data as
    downloaded: the monthly annual rate of change of the HICP (dataset
    \texttt{prc\_hicp\_manr}, an Excel extract in \texttt{ecb/}) and the chain-linked
    volume index of euro-area GDP, seasonally and calendar adjusted (dataset
    \texttt{namq\_10\_gdp}, file \texttt{ecb/namq\_10\_gdp\_EA\_RGDP.xml}). Forecasters
    need at least $15$ observations.
  \item \emph{Good Judgment Project} (\S\ref{sec:apps-gjp}). The question file
    \texttt{gjp/ifps.csv} and the individual forecasts of tournament years~1 and~2,
    \texttt{gjp/survey\_fcasts.yr1.csv} and \texttt{gjp/survey\_fcasts.yr2.csv}, come
    from the Harvard Dataverse release \citep{gjpdata2016}, licensed CC0.
    \texttt{data.gjp.brier\_panel} keeps closed questions with a resolved outcome,
    drops withdrawn forecasts and orders the questions by close date.
  \item \emph{M3} (\S\ref{sec:apps-mcomp}). The series, the held-out values and the
    24 submitted forecasts are the data sets \texttt{M3} and \texttt{M3Forecast} of
    the R package \texttt{Mcomp}, version~2.8 \citep{hyndman2018mcomp}. Unlike the
    other sources, these need no manual download.
    \texttt{data.mcomp.load\_m3\_official} fetches the package archive from CRAN,
    checks its SHA-256 checksum, reads the two R data files with the Python package
    \texttt{rdata}, and caches them as one long table,
    \texttt{mcomp/m3\_official\_long.csv}, with one row per series and horizon and
    one column per method. \texttt{data.mcomp.official\_panel} builds the Paris panel
    from the rows whose description is ``Average temperature in Paris'' (series
    N2784--N2800), keeping the blocks in calendar order.
\end{itemize}
The Philadelphia Fed and Eurostat files grow as new data are released, and older
observations are occasionally corrected. A later download therefore adds observations
at the end of each panel and can change which forecasters are selected, so results
computed from it will differ from those reported. The panels used here end
with the 2025Q1 survey (U.S.\ output), the 2025Q3 survey (U.S.\ unemployment) and
target year 2025 (euro area). The M3 data cannot drift in this way, because the
package version is fixed and its checksum is verified.

\section{Software and tests}
\label{app:impl-env}

The results were computed with the software in Table~\ref{tab:impl-env}. A fresh
environment is set up from the project directory with
\begin{verbatim}
python3.11 -m venv .venv
.venv/bin/pip install -e .
.venv/bin/pip install scipy matplotlib openpyxl rdata jupyter pytest
\end{verbatim}
pinning the versions of Table~\ref{tab:impl-env} to match the reported numbers
exactly. The first run of the M3 loader needs an internet connection.

\tabImplEnv

\subsection{Unit tests}
\label{app:impl-tests}

The directory \texttt{tests/} contains 65 unit tests, run with
\texttt{python -m pytest}. All pass in about three seconds. One of them
downloads the M3 data from CRAN and runs only when the environment variable
\texttt{RANKCI\_NETWORK} is set; the others run offline.
Table~\ref{tab:impl-tests} lists what they check. The estimator tests target the
properties the main text relies on: the bandwidth formula \eqref{eq:andrews}, the
diagonal-consistency and nesting of $\Omegahat$ (\S\ref{sec:omega-nesting}), the PSD
statements of \S\ref{sec:caveat}, the agreement of the two constructions on a balanced
panel (\S\ref{sec:equivalence}), the cancellation of the common shock in the
simulated panels, the block-bootstrap index construction and the equivalence of the
bootstrap variance with the Newey--West variance (\S\ref{par:mbb}), and the coverage
argument for the $\tau$-best projection (\S\ref{sec:taubest}). The marginal rank sets and
the bootstrap variant of the direct $\tau$-best procedure are covered by a smoke test
only, which checks that they run on an unbalanced panel and return well-formed sets. The
Gaussian variant of the direct $\tau$-best procedure is exercised only through the
notebooks and scripts.

\tabImplTests

\section{Reproducing the results}
\label{app:impl-repro}

With the data in place and the package installed, the results are reproduced in three
steps.
\begin{enumerate}[nosep,leftmargin=*]
  \item Run the tests: \texttt{python -m pytest}.
  \item Execute the notebooks of Table~\ref{tab:impl-repro}, for example
    \texttt{jupyter nbconvert --to notebook --execute --inplace
    notebooks/01\_toydataset.ipynb}. Each notebook reads the data by paths relative
    to its own folder, which is where \texttt{nbconvert} runs it.
    \texttt{02\_primary\_analysis.ipynb} writes its tables and figures directly into
    \texttt{thesis/tables/} and \texttt{thesis/figures/}.
  \item Run the two scripts, \texttt{python scripts/make\_thesis\_figures.py} and
    \texttt{python scripts/make\_thesis\_tables.py}. They write into
    \texttt{thesis/figures/} and \texttt{thesis/tables/}.
\end{enumerate}
Table~\ref{tab:impl-repro} lists the file and the settings behind each result. As a
check, the crossover sweep of Table~\ref{tab:mc-crossover}, both columns of
Table~\ref{tab:mcomp}, and Tables~\ref{tab:spf-indicators}, \ref{tab:ngdp}
and~\ref{tab:threeregimes} were rerun from scratch with these settings, and all
reproduce every entry. On the machine of Table~\ref{tab:impl-env} the crossover sweep
takes about twelve seconds. The table script takes about a minute with
\texttt{--skip-mc}; the Monte Carlo tables of Appendix~\ref{app:additional} add about
five minutes. On the M3 panel each simulation column takes about a second.

\tabImplRepro

\section{Known limitations of the implementation}
\label{app:impl-limits}

\begin{itemize}[leftmargin=*]
  \item \emph{The direct $\tau$-best procedure} is implemented with MDS draws but not
    with $\Omegahat$. The $\tau$-best set on M3 (\S\ref{sec:apps-mcomp-results}) uses
    the MDS procedure with $B=20000$; it equals the projection of the direct
    construction's rank sets. The $\tau$-best sets reported for NGDP and GJP are
    projections of the direct construction's joint rank sets (\S\ref{sec:taubest}).
  \item \emph{The two lag conventions} of \S\ref{app:impl-details} differ on overlaps
    with gaps. Each route is internally consistent, but on such panels the MDS and
    direct routes do not studentize by exactly the same standard errors. The block
    bootstrap resamples blocks of consecutive rows, counting lags in rows as the direct
    estimator does, but takes its block length from the bandwidths of the concatenated
    overlap series.
  \item \emph{Cost.} The direct estimator assembles $q(q+1)/2$ entries, each a
    weighted sum over up to $L_{ab}+1$ lags, in plain Python loops; with $p=24$
    ($q=276$) this takes about a second.
\end{itemize}

\chapter{Additional Results}
\label{app:additional}

The tables below support statements that Chapters~\ref{ch:mc} and~\ref{ch:apps} make
in the text only. Section~\ref{app:add-mc} completes the Monte Carlo study.
Sections~\ref{app:add-macro} and~\ref{app:add-gjp} list every rank set of the
macroeconomic panels and of the tournament. Section~\ref{app:add-m3} gives the
block-level losses behind the M3 results and two robustness checks. As in the main
text, every rank set comes from the simulation route; the block bootstrap appears only
as the benchmark column of the Monte Carlo tables. \texttt{scripts/make\_thesis\_tables.py}
writes every table here (Table~\ref{tab:impl-repro}).

\section{Monte Carlo: the full design grid}
\label{app:add-mc}

Table~\ref{tab:mc-coverage} reports eight of the eleven designs of the Monte Carlo study.
Tables~\ref{tab:app-mc-coverage} and~\ref{tab:app-mc-width} report all eleven, with the
same designs and seeds, computed with the software of Table~\ref{tab:impl-env}. The three
designs added here are white-noise idiosyncratic errors ($\rho=0$), a milder imbalance
of $0.2$, and the tie design of \S\ref{sec:mc-coverage}, $\theta=(0,0,0.8,0.8,1.6)$. All
other parameters are those of the baseline: $p=5$, $T=200$, $\theta$ equally spaced on
$[0,1.6]$, $\rho=0.3$ and $\sigma_c=1$. The non-stationary design has five volatility
episodes of scale $12$, and the realistic design combines them with an imbalance of
$0.4$.

\paragraph{Coverage.} Joint coverage is one or nearly one ($\ge 0.998$) on every design
with distinct means except $\rho=0.9$ ($0.966$--$0.968$), well above the nominal $0.95$
(Table~\ref{tab:app-mc-coverage}). This over-coverage is a property of rank sets, not a
sign of an inflated critical value. A rank set misses the
true rank only if the procedure confirms an ordering in the wrong direction. With
adjacent means $0.4$ apart such an error is rare, and the level $\alpha$ is approached
only for a pair with equal means. The tie design has two such pairs. A tied pair is
confidently ordered in $2.0\%$ of its replications under MDS and the block bootstrap
and in $1.9\%$ under the direct construction, below the nominal $5\%$. The smallest
marginal coverage, that of the least-well-covered forecaster's own set, is at least
$0.982$ and equals the joint coverage on every design except $\rho=0.9$. The lowest
coverage occurs at $\rho=0.9$: $0.968$ for both simulation routes and $0.966$ for the
block bootstrap. With $T=200$ and strong persistence, kernel
estimators of the long-run variance are biased downward \citep{andrews1991}, so the
studentized contrasts are more dispersed than the unit-variance draws they are compared
with. The block bootstrap inherits the same bias, because its block length reproduces
the Bartlett weights of these estimators (\S\ref{par:mbb}).

\paragraph{Width, projection and stepdown.} Table~\ref{tab:app-mc-width} gives the
widths. Without serial dependence ($\rho=0$) the sets are sharpest, with a mean width of
about $0.11$ under all three routes. The milder imbalance of $0.2$ behaves like the
imbalance of $0.4$ on a smaller scale: MDS ($0.662$) is marginally sharper than the
direct construction ($0.678$) and the block bootstrap ($0.682$), which are
indistinguishable. The projection distance of $\Omegahat$ is at most $0.004$ on the balanced stationary
designs, essentially zero under imbalance alone, and larger only under
non-stationarity ($0.034$, and $0.067$ on the realistic design): with a separate
bandwidth for each difference series the estimator is not guaranteed to be positive
semi-definite, and heterogeneous volatility makes the bandwidths most heterogeneous
(\S\ref{sec:caveat}). The MDS matrix is essentially positive semi-definite (distance at most $0.001$) except
on the realistic design.
A positive distance does not by itself cost the direct construction anything, because
its draws are studentized by the diagonal of the projected matrix. The stepdown takes
about two rounds on the stationary designs ($2.0$--$2.2$ on average) and fewer on the
non-stationary ones ($1.4$--$1.8$), where the sets are wide and fewer pairs are
confirmed.

\paragraph{Crossover.} Table~\ref{tab:app-mc-crossover} completes
Table~\ref{tab:mc-crossover}. Both routes cover in every replication at every crisis
scale. The projection distance of $\Omegahat$ grows from essentially zero to $0.29$ and
that of the MDS matrix from $0.001$ to $1.0$, and the direct construction is the
sharper of the two from scale $6$ on. The two distances refer to different matrices, the
$q\times q$ covariance of the differences and the $p\times p$ level covariance, so their
sizes are not directly comparable.
What matters is whether the draws are studentized consistently with the projected
matrix. That holds for the direct route and fails for the MDS route
(\S\ref{sec:mc-mechanism}).

\tabAppMcCoverage
\tabAppMcWidth
\tabAppMcCrossover

\section{Macroeconomic panels: all rank sets}
\label{app:add-macro}

Table~\ref{tab:spf-indicators} summarizes each U.S.\ panel by the number of separated
forecasters and the set of the forecaster with the smallest mean loss.
Table~\ref{tab:app-spf-sets} lists all twenty-four sets under both simulation routes.
\begin{itemize}[leftmargin=*]
  \item The MDS route returns the trivial $[1,8]$ for every forecaster on all three
    panels. These panels are as unbalanced as the five-forecaster NGDP panel of
    Table~\ref{tab:ngdp} (each forecaster reports in about half of the quarters), and
    the MDS route loses calibration in the same way.
  \item A small mean loss does not by itself give a good rank set. On NGDP, forecaster~40
    has the smallest mean loss but a trivial set. The only set with a lower bound above
    one belongs to forecaster~65, whose mean loss is the second smallest. On
    unemployment, forecaster~65 again has the second smallest mean loss but is confined
    to the last three ranks. The contrasts compare forecasters on their shared quarters
    (\S\ref{sec:setup}), whereas the mean loss also reflects which quarters a forecaster
    happened to cover.
  \item The best forecaster on unemployment, forecaster~84, has the set $[1,6]$. It is
    confirmed better than two rivals and no rival is confirmed better than it, but the
    other five cannot be ordered against it, so it cannot be named best.
\end{itemize}

Table~\ref{tab:app-ecb-sets} gives the euro-area sets. The two non-trivial HICP sets are
those of forecasters~16 ($[1,7]$) and~89 ($[2,8]$). On real GDP growth every set is
trivial. The two routes agree exactly on both panels. The panels are almost balanced
(every forecaster covers 26 or 27 of the 27 target years), so the equivalence of
\S\ref{sec:equivalence} applies.

\tabAppSpfSets
\tabAppEcbSets

\section{The geopolitical tournament: all rank sets}
\label{app:add-gjp}

Table~\ref{tab:app-gjp} lists the twenty sets plotted in
Figure~\ref{fig:gjp-caterpillar}. Three forecasters have a lower bound of one
(803, 1038 and 1000); they form the $\tau$-best set for $\tau=1$ of
\S\ref{sec:apps-gjp}. Forecaster~660 is pinned to the last rank, so no other forecaster
can rank twentieth and every other upper bound is at most $19$. The two routes differ
by one rank for two forecasters and agree for the rest. The panel behind them is almost
balanced: each forecaster answered at least 210 of the 212 questions.

The last two columns support the split-half check of \S\ref{sec:apps-gjp}. The questions,
ordered by closing date, are split into the odd- and the even-numbered ones, and the
forecasters are ranked by mean Brier score on each half. Forecaster~803 is first and
forecaster~660 last in both halves. In the middle the two rankings differ by up to eight
places (forecaster~3631 is fifth on one half and thirteenth on the other), which matches
the wide sets in the middle of the table.

\tabAppGjp

\section{M3: block-level losses and robustness}
\label{app:add-m3}

\subsection{Losses by block}
\label{app:add-m3-blocks}

Table~\ref{tab:app-m3-blocks} gives the sMAPE of six submissions in each block of the
Paris record: B-J automatic, which has the smallest mean sMAPE; Theta, the overall
winner of M3; PP-Autocast; and the three methods pinned to the bottom ranks.
\begin{itemize}[leftmargin=*]
  \item The leading methods move together. Their sMAPE lies between about $6.5$ and $19$
    depending on the block, and the hard blocks are hard for all of them. This
    co-movement is the common shock that cancels in the loss differences.
  \item Each Flores--Pearce variant forecasts a constant value in twelve of the
    seventeen blocks (marked with a star), with an sMAPE of $28$--$50$ in those blocks.
    In the five blocks in which Flores--Pearce2 does not forecast a constant, it is as
    accurate as the leaders. Flores--Pearce1 is far behind them in every block.
  \item Robust-Trend has an sMAPE above $25$ in eight blocks and is close to the leaders
    in the other nine.
  \item PP-Autocast is close to the leaders in most blocks; the exception is 1953--1960, where its
    sMAPE is $37.0$ against about $8$. This single block drives the variance of its loss
    differences and hence the width of its set, $[1,21]$.
\end{itemize}

\tabAppMThreeBlocks

\subsection{Robustness of the rank sets}
\label{app:add-m3-robust}

Table~\ref{tab:app-m3-robust} gives the sets behind the two robustness checks of
\S\ref{sec:apps-mcomp-results}. At $\alpha=0.05$ the same twelve methods keep rank one
and the same three remain the only methods excluded from the top twenty, although
Robust-Trend's set widens to $[21,24]$ and PP-Autocast's to $[1,22]$. Most sets widen by at most one rank. The
largest changes are the lower bounds of AutoBox3 (from $7$ to $3$) and ARARMA (from $13$
to $9$). On the panel of block averages ($n=17$) fourteen methods keep rank one: Naive2
and AutoBox1 join the twelve. The sets are not uniformly wider, however. AAM2 and
AutoBox3, for example, gain a lower bound of $9$. Averaging over the eighteen horizons
removes the dependence between horizons within a block, and with only seventeen
observations the standard errors are estimated far less precisely, so the two panels
need not give nested sets.

\tabAppMThreeRobust

\chapter{Dependence in the SPF Forecast Errors}
\label{app:spf-dependence}

Here I document the evidence summarized in \S\ref{sec:apps-spf-dependence}.
\texttt{notebooks/02\_primary\_analysis.ipynb} produces all numbers, tables and figures.

\section{Error panels}
\label{app:spf-dep-panels}

\paragraph{U.S.\ SPF.} For indicator $v\in\{\text{NGDP},\text{RGDP},\text{UNEMP}\}$,
forecaster $j$ and survey quarter $t$, the SPF microdata give the forecast of the quarter
after the survey (SPF horizon index $3$). The realization is the advance estimate: the
value of the target quarter in the first vintage of the Philadelphia Fed Real-Time Data
Set that contains it. Monthly unemployment is averaged to quarters. NGDP and RGDP are
forecast and realized in levels. The error is $e_{t,j}=\text{forecast}-\text{realization}$,
indexed by survey quarter.

\paragraph{Euro-area SPF.} For HICP inflation and real GDP growth, I use the forecast of
the current calendar year's average year-on-year rate made in the first-quarter survey
round, the panel ranked in \S\ref{sec:apps-ecb}. Realizations are taken from
Eurostat as downloaded, not from real-time vintages. The HICP realization is the mean of
the twelve monthly year-on-year rates. The real-GDP realization is the mean of the four
quarterly year-on-year growth rates of the chain-linked volume index. Errors are indexed
by target year.

\paragraph{Construction.} Each panel is placed on a regular time grid, so that a missing
survey is a missing row and lag $k$ always means $k$ periods apart. Per-forecaster
statistics use every forecaster with at least $20$ (U.S.) or $15$ (euro area)
observations. Statistics on loss differences use the eight best-covered forecasters, the
same selection as the rank inference. Table~\ref{tab:spf-dep-panels} describes the panels,
and Figure~\ref{fig:spf-dep-errors} plots the errors.

\tabSpfDepPanels
\figSpfDepErrors

\section{Serial dependence}
\label{app:spf-dep-serial}

Because of the gaps, autocorrelations are computed from all pairs of periods $k$ apart in
which the forecaster reported. Let $\mathcal{P}_k=\{t: e_{t,j}\text{ and }e_{t+k,j}\text{
observed}\}$ and $n_k=|\mathcal{P}_k|$. Then
\begin{equation}
  \hat\rho_k \;=\; \frac{n_k^{-1}\sum_{t\in\mathcal{P}_k}(e_{t,j}-\bar e_j)(e_{t+k,j}-\bar e_j)}
                     {n^{-1}\sum_{t}(e_{t,j}-\bar e_j)^2},
\end{equation}
reported only when $n_k\ge 8$. Under no serial correlation, $\hat\rho_k$ is approximately
$\Normal(0,1/n_k)$, which gives the band $\pm1.96/\sqrt{n_k}$. Joint whiteness up to lag
$K$ is tested by the Box--Pierce statistic $Q=\sum_{k=1}^K n_k\hat\rho_k^2$, compared with
$\chi^2_K$ \citep{boxpierce1970}. I use $K=4$ (one year) for the U.S.\ and $K=2$ for the
euro area, and compute the same statistics on the losses $e_{t,j}^2$. For the loss
differences $\dser{jk}_t$, I report the AR(1) coefficient and the Andrews bandwidth that
the HAC layer actually uses (\S\ref{sec:bandwidth}). Both are computed on the pair's
overlap, as in the estimator.

Some serial correlation is expected even from efficient forecasters. An optimal $h$-step
forecast error is at most $(h-1)$-dependent \citep{dieboldmariano1995}. The SPF forecast
for the quarter after the survey is made before the survey quarter itself is published,
so it is effectively a two-step forecast, and the errors of consecutive surveys share one
unobserved quarter. This shared quarter gives an MA(1) structure. The euro-area current-year targets do not
overlap from one year to the next.

Table~\ref{tab:spf-dep-serial} and Figure~\ref{fig:spf-dep-acf} report the results. The
U.S.\ errors and losses have a clear first-order autocorrelation and little beyond it. The
euro-area panels show none, but with $27$ annual observations the white-noise band is
$\pm0.38$, so only strong persistence could be detected. The loss differences are only
mildly persistent, and the plug-in bandwidth is correspondingly short.

\tabSpfDepSerial
\figSpfDepAcf

\section{Cross-sectional dependence}
\label{app:spf-dep-cross}

For each pair of forecasters, $\hat\rho_{jk}$ is the Pearson correlation of their errors
over the $\Tjk$ common periods, required to number at least eight. I summarize the
pairwise correlations in three ways.
\begin{itemize}[leftmargin=*]
  \item \emph{Mean correlation} $\bar\rho$ over all pairs, for errors and for losses, and
    its Spearman analogue $\bar\rho_S$, which is insensitive to the size of extreme
    periods.
  \item \emph{Pesaran's CD statistic} for unbalanced panels \citep{pesaran2021},
    \begin{equation}
      \mathrm{CD} \;=\; \frac{1}{\sqrt{P}}\sum_{j<k}\sqrt{\Tjk}\,\hat\rho_{jk},
    \end{equation}
    over the $P$ valid pairs. It is approximately $\Normal(0,1)$ under cross-sectional
    independence.
  \item \emph{The period-effect share} $R^2=1-\sum_{t,j}(e_{t,j}-\bar e_t)^2/
    \sum_{t,j}(e_{t,j}-\bar e)^2$. Here $\bar e_t$ is the cross-sectional mean of the
    forecasters reporting in period $t$, and only periods with at least two reports are
    used. With no common component, $R^2$ is about the number of periods divided by the
    number of observations, reported in parentheses.
\end{itemize}
Table~\ref{tab:spf-dep-cross} reports the results. Every panel rejects cross-sectional
independence by a wide margin. The Spearman correlations are only slightly lower than
the Pearson ones, and dropping 2020--21 changes the mean error correlation by at most
$0.06$ (from $0.76$--$0.97$ to $0.73$--$0.92$). The dependence is therefore a persistent
feature of the data, not the product of one crisis.

\tabSpfDepCross

\section{The common shock in the loss differences}
\label{app:spf-dep-diff}

The decomposition \eqref{eq:decomp} treats the common shock as additive in the loss, and
it then cancels exactly in $\dser{jk}_t$. For squared errors the common component enters
the error instead. With $e_{t,j}=f_t+u_{t,j}$,
\begin{equation}
  \dser{jk}_t \;=\; (u_{t,j}^2-u_{t,k}^2) \;+\; 2f_t\,(u_{t,j}-u_{t,k}).
  \label{eq:reentry}
\end{equation}
The common error cancels from the level of $\dser{jk}_t$ but not from its scale. If the
idiosyncratic errors are symmetric and independent of $f_t$,
\begin{equation}
  \Var\big(\dser{jk}_t\mid f_t\big) \;=\; \Var\big(u_{t,j}^2-u_{t,k}^2\big)
  \;+\; 4f_t^2\,\Var\big(u_{t,j}-u_{t,k}\big),
\end{equation}
which grows with the squared common error. Table~\ref{tab:spf-dep-diff} checks both
predictions.

Differencing removes most of the loss variance: $82\%$ to essentially $100\%$ on the
reported panels. What remains moves with the common shock. The squared loss differences
are rank-correlated with the squared consensus error at $0.63$--$0.89$ (and at
$0.61$--$0.87$ without 2020--21). Between $94\%$ and $100\%$ of their sum falls in the
tenth of periods with the largest common error ($73$--$99\%$ without 2020--21). The
pairwise difference series are therefore heteroskedastic, with variance concentrated in
a few episodes. On an unbalanced panel, different pairs contain different
such episodes, which is the source of the inconsistency among pairwise standard errors
analysed in \S\ref{sec:mc-mechanism}.

\tabSpfDepDiff

\section{Level forecasts, base-year changes and a growth-rate check}
\label{app:spf-dep-rebasing}

The U.S.\ NGDP and RGDP panels are in levels. When a benchmark revision or base-year
change is published between the survey and the advance release of its target, the
forecast and the realization are on different bases. Every forecaster's error then
shifts by the same amount, typically in two consecutive surveys. These shifts are the
recurrent spikes of the U.S.\ RGDP panel in Figure~\ref{fig:spf-dep-errors}. They
inflate both the cross-sectional correlation and the lag-1 autocorrelation, and they add
episodes of large variance to the loss differences.

As a check, I re-express the same forecasts as annualized quarter-on-quarter growth
rates. The forecast growth rate is computed from the forecaster's own submitted path, and
the realized growth rate from the advance vintage of the target. Both levels in each
ratio are therefore on the same base. These are the rows marked $^\dagger$ in
Tables~\ref{tab:spf-dep-panels}--\ref{tab:spf-dep-diff}.
\begin{itemize}[leftmargin=*]
  \item The rebasing spikes disappear; the only extreme periods left are 2020Q1--Q2.
  \item Cross-sectional dependence remains strong: the mean error correlation is
    $0.74$--$0.77$ and the CD statistic exceeds $240$.
  \item Serial correlation is weaker (median lag-1 autocorrelation about $0.1$) but still
    flags $29\%$ of forecasters, far above the $5\%$ expected by chance.
  \item The loss differences remain heteroskedastic in the common shock.
\end{itemize}
The conclusions of \S\ref{sec:apps-spf-dependence} therefore do not depend on the level
artefacts.

\end{appendix}
